% Part: normal-modal-logic
% Chapter: axioms-systems
% Section: logics-proofs

\documentclass[../../../include/open-logic-section]{subfiles}

\begin{document}

\olfileid[id]{nml}{prf}{prf}

\olsection{\usetoken{P}{derivation} dan Sistem Modal}

Mula-mula kita mendefinisikan !!a{derivation} bagi logika modal normal. Secara
kasar, !!a{derivation} merupakan barisan !!{formula}s dengan setiap
!!{element} merupakan suatu instansiasi substitusi dari salah satu dari
sejumlah \emph{aksioma}, atau mengikuti dari !!{element}s sebelumnya melalui
salah satu dari beberapa aturan inferensi. Bagi logika modal normal, semua
instansiasi tautologi\iftag{prvDiamond}{, \Ax{K}, dan~\Dual{}}{
dan~\Ax{K}} dihitung sebagai aksioma. Hal ini menghasilkan sistem
modal~$\Log{K}$, logika modal normal terkecil. Meskipun demikian, kita dapat
menambahkan aksioma lain untuk memperoleh sistem-sistem lain. Aturannya selalu
modus ponens~\MP{} dan nesesitasi~\Nec.

\begin{defn}
  Diberikan sistem modal $\Log{K} !A_1 \dots !A_n$ dan !!a{formula}~$!B$,
  kita mengatakan bahwa $!B$ \emph{!!{derivable}} dalam
  $\Log{K} !A_1 \dots !A_n$, ditulis
  $\Log{K} !A_1 \dots !A_n\Proves !B$, jika dan hanya jika terdapat
  !!{formula}s $!C_1$, \dots, $!C_k$ sedemikian sehingga $!C_k=!B$ dan setiap
  $!C_i$ merupakan instansiasi tautologis, instansiasi salah satu dari
  $\Ax{K}$,\iftag{prvDiamond}{ $\Dual$,}{} $!A_1$, \dots, $!A_n$, atau
  mengikuti dari !!{formula}s sebelumnya melalui aturan~\MP{} atau~\Nec.
\end{defn}

Proposisi berikut memungkinkan kita menunjukkan bahwa $!B\in\Sigma$ dengan
memperlihatkan suatu $\Sigma$-!!{derivation} dari~$!B$.

\begin{prop}
  $\Log{K} !A_1 \dots !A_n =
  \Setabs{!B}{\Log{K} !A_1 \dots !A_n \Proves !B}$.
\end{prop}

\begin{proof}
  Kita menggunakan induksi pada panjang !!{derivation}s untuk menunjukkan
  bahwa
  $\Setabs{!B}{\Log{K} !A_1 \dots !A_n\Proves !B}\subseteq
  \Log{K} !A_1 \dots !A_n$.

  Jika !!{derivation} dari~$!B$ mempunyai panjang~$1$, derivasi itu memuat satu
  !!{formula}. !!^{formula} tersebut tidak mungkin mengikuti dari formula
  sebelumnya melalui~\MP{} atau~\Nec, sehingga harus merupakan instansiasi
  tautologis, instansiasi~\Ax{K},\iftag{prvDiamond}{ $\Dual$,}{} atau
  instansiasi salah satu dari $!A_1$, \dots,~$!A_n$. Akan tetapi,
  $\Log{K}!A_1\dots!A_n$ juga memuat semua formula tersebut, sehingga
  $!B\in\Log{K}!A_1\dots !A_n$.

  Jika !!{derivation} dari~$!B$ mempunyai panjang~$>1$, $!B$ juga mungkin
  diperoleh melalui~\MP{} atau~\Nec{} dari !!{formula}s yang tidak muncul
  sebagai baris terakhir !!{derivation}. Jika $!B$ mengikuti dari $!C$ dan
  $!C\lif !B$ melalui~\MP, maka menurut hipotesis induksi
  $!C$ dan $!C\lif !B\in\Log{K}!A_1\dots !A_n$. Akan tetapi, setiap logika
  modal tertutup di bawah modus ponens, sehingga
  $!B\in\Log{K}!A_1\dots !A_n$. Jika $!B\equiv\Box !C$ mengikuti dari~$!C$
  melalui~\Nec, maka menurut hipotesis induksi
  $!C\in\Log{K}!A_1\dots !A_n$. Akan tetapi, setiap logika modal normal
  tertutup di bawah~$\Nec$, sehingga $!B\in\Log{K}!A_1\dots!A_n$.

  Inklusi konvers diperoleh dengan menunjukkan bahwa
  $\Sigma=\Setabs{!B}{\Log{K}!A_1\dots!A_n\Proves !B}$ merupakan logika modal
  normal yang memuat semua instansiasi $!A_1$, \dots, $!A_n$, lalu menggunakan
  fakta bahwa menurut definisi $\Log{K}!A_1\dots!A_n$ merupakan logika
  terkecil semacam itu.
  \begin{enumerate}
    \item Setiap tautologi~$!B$ merupakan instansiasi tautologis, sehingga
      $\Log{K}!A_1\dots!A_n\Proves !B$; jadi, $\Sigma$ memuat semua tautologi.
    \item Jika $\Log{K}!A_1\dots!A_n\Proves !C$ dan
      $\Log{K}!A_1\dots!A_n\Proves !C\lif !B$, maka
      $\Log{K}!A_1\dots!A_n\Proves !B$: gabungkan !!{derivation} dari~$!C$
      dengan derivasi dari~$!C\lif !B$, lalu tambahkan baris~$!B$. Baris
      terakhir dijustifikasi oleh~\MP{}. Jadi, $\Sigma$ tertutup di bawah
      modus ponens.
    \item Jika $!B$ mempunyai !!a{derivation}, setiap instansiasi substitusi
      dari~$!B$ juga mempunyai !!a{derivation}: terapkan substitusi tersebut
      pada setiap !!{formula} dalam !!{derivation}. (Latihan: buktikan melalui
      induksi pada panjang !!{derivation}s bahwa hasilnya juga merupakan
      !!a{derivation} yang benar.) Jadi, $\Sigma$ tertutup di bawah substitusi
      seragam. (Kini kita telah menetapkan bahwa $\Sigma$ memenuhi semua
      kondisi logika modal.)
    \item Kita mempunyai $\Log{K}!A_1\dots!A_n\Proves\Ax{K}$, sehingga
      $\Ax{K}\in\Sigma$.
    \tagitem{prvDiamond}{Kita mempunyai
      $\Log{K}!A_1\dots!A_n\Proves\Dual$, sehingga $\Dual\in\Sigma$.}{}
    \item Jika $\Log{K}!A_1\dots!A_n\Proves !C$, baris tambahan~$\Box !C$
      dijustifikasi oleh~\Nec. Akibatnya, $\Sigma$ tertutup di bawah~\Nec.
      Jadi, $\Sigma$ normal.
  \end{enumerate}
\end{proof}

\end{document}
